\chapter{Types of Recursion}

Structural recursion is the default, but some problems need more. CS 135
classifies recursive functions as follows:

\begin{itemize}
  \item \textbf{Pure structural}: recursive calls only on the self-referential parts; other parameters unchanged.
  \item \textbf{Accumulative} (structural with an accumulator): recursion is structural on one parameter, but extra parameters \emph{accumulate} information about what has been seen so far.
  \item \textbf{Mutual}: two or more functions call each other.
  \item \textbf{Generative}: recursive calls are on values \emph{computed} from the input, not just its parts (Chapter~\ref{ch:generative}).
\end{itemize}

\section{Accumulative recursion}

\subsection{Motivation: reversing a list}
A structural solution:
\begin{lstlisting}[style=racket]
;; (my-reverse lst) produces lst in reverse order.
;; my-reverse: (listof Any) -> (listof Any)
(define (my-reverse lst)
  (cond [(empty? lst) empty]
        [else (append (my-reverse (rest lst)) (list (first lst)))]))
\end{lstlisting}
Each \code{append} walks the entire reversed rest, so the total work is
$1 + 2 + \cdots + n = \Theta(n^2)$. With an accumulator we can do it in linear
time:

\begin{lstlisting}[style=racket]
;; (my-reverse lst) produces lst in reverse order.
;; Examples:
(check-expect (my-reverse '(1 2 3)) '(3 2 1))
(check-expect (my-reverse empty) empty)

;; my-reverse: (listof Any) -> (listof Any)
(define (my-reverse lst) (rev/acc lst empty))

;; (rev/acc lst acc) produces the elements of lst in reverse order
;;   followed by acc.
;; rev/acc: (listof Any) (listof Any) -> (listof Any)
(define (rev/acc lst acc)
  (cond [(empty? lst) acc]
        [else (rev/acc (rest lst) (cons (first lst) acc))]))
\end{lstlisting}

\paragraph{Trace.}
\begin{lstlisting}[style=racket]
(my-reverse (list 1 2 3))
=> (rev/acc (list 1 2 3) empty)
=> (rev/acc (list 2 3) (list 1))
=> (rev/acc (list 3) (list 2 1))
=> (rev/acc empty (list 3 2 1))
=> (list 3 2 1)
\end{lstlisting}

\begin{important}[Accumulative recursion recipe]
\begin{itemize}
  \item A \textbf{wrapper} function calls the helper with the \emph{initial} accumulator value.
  \item The helper's purpose explains what the accumulator \emph{means} (the accumulator invariant).
  \item In the recursive call the structural parameter moves to the rest, and the accumulator is updated to reflect the element just processed.
  \item In the base case the answer is usually (computed from) the accumulator.
\end{itemize}
\end{important}

\subsection{Template}
\begin{lstlisting}[style=racket]
;; acc-template: (listof X) Y -> Any
(define (acc-template lst acc)
  (cond [(empty? lst) (... acc ...)]
        [else (acc-template (rest lst)
                            (... (first lst) ... acc ...))]))
\end{lstlisting}

\subsection{Maximum with an accumulator}
\begin{lstlisting}[style=racket]
;; (list-max lon) produces the largest element of lon.
;; list-max: (ne-listof Num) -> Num
(check-expect (list-max '(1 7 3)) 7)
(define (list-max lon) (max/acc (rest lon) (first lon)))

;; (max/acc lon best) produces the max of best and all elements of lon.
;; max/acc: (listof Num) Num -> Num
(define (max/acc lon best)
  (cond [(empty? lon) best]
        [(> (first lon) best) (max/acc (rest lon) (first lon))]
        [else (max/acc (rest lon) best)]))
\end{lstlisting}

\subsection{Average in one pass}
Two accumulators can track several quantities at once:
\begin{lstlisting}[style=racket]
;; (average lon) produces the mean of lon.
;; average: (ne-listof Num) -> Num
(check-expect (average '(1 2 3 6)) 3)
(define (average lon) (avg/acc lon 0 0))

;; (avg/acc lon total count): total and count are the sum and number
;;   of elements already processed.
;; avg/acc: (listof Num) Num Nat -> Num
(define (avg/acc lon total count)
  (cond [(empty? lon) (/ total count)]
        [else (avg/acc (rest lon) (+ total (first lon)) (add1 count))]))
\end{lstlisting}

\begin{aside}
Accumulative functions are often \emph{tail recursive}: the recursive call is
the whole answer, so nothing waits to be done after it returns. In the stepper
the expression stays the same size, unlike \code{(+ 2 (+ 3 (sum ...)))}, which
grows. This matters for space efficiency in real implementations.
\end{aside}

\subsection{Context-dependent processing}
Accumulators also give the helper access to context such as the position of an
element or the previous element:
\begin{lstlisting}[style=racket]
;; (increasing? lon) produces true if each element is larger than the
;;   one before it.
;; increasing?: (listof Num) -> Bool
(check-expect (increasing? '(1 3 8)) true)
(check-expect (increasing? '(1 3 3)) false)
(define (increasing? lon)
  (or (empty? lon) (inc/acc (rest lon) (first lon))))

;; inc/acc: (listof Num) Num -> Bool
(define (inc/acc lon prev)
  (cond [(empty? lon) true]
        [else (and (> (first lon) prev)
                   (inc/acc (rest lon) (first lon)))]))
\end{lstlisting}

\section{Mutual recursion}

Mutual recursion arises naturally from mutually recursive data definitions
(Chapter~\ref{ch:trees}), but even simple examples exist:

\begin{lstlisting}[style=racket]
;; (is-even? n) produces true if n is even.
;; is-even?: Nat -> Bool
(define (is-even? n)
  (cond [(zero? n) true]
        [else (is-odd? (sub1 n))]))

;; (is-odd? n) produces true if n is odd.
;; is-odd?: Nat -> Bool
(define (is-odd? n)
  (cond [(zero? n) false]
        [else (is-even? (sub1 n))]))
\end{lstlisting}

\begin{example}[Alternating elements]
Keep the elements in odd positions (1st, 3rd, \ldots):
\begin{lstlisting}[style=racket]
;; keep-alternates: (listof Any) -> (listof Any)
(check-expect (keep-alternates '(a b c d e)) '(a c e))
(define (keep-alternates lst)
  (cond [(empty? lst) empty]
        [else (cons (first lst) (skip-one (rest lst)))]))

;; skip-one: (listof Any) -> (listof Any)
(define (skip-one lst)
  (cond [(empty? lst) empty]
        [else (keep-alternates (rest lst))]))
\end{lstlisting}
Each function corresponds to a ``state'': keeping or skipping.
\end{example}

\section{Generative recursion (preview)}
In generative recursion the argument to the recursive call is generated by
some computation, e.g.\ \code{(gcd m (remainder n m))}. Termination is no
longer automatic and requires an argument. Chapter~\ref{ch:generative} covers
this in depth.

\section{Choosing a kind of recursion}
\begin{itemize}
  \item Start with structural recursion; it is simplest and terminates automatically.
  \item Add an accumulator when the answer depends on what came \emph{before} (context), or to avoid repeated work like \code{append} in reverse.
  \item Use mutual recursion when the data definition is mutually recursive or the problem has distinct states.
  \item Use generative recursion only when the problem's structure, not the data's, dictates the decomposition.
\end{itemize}

\section{Summary}
\begin{itemize}
  \item Accumulators carry information about the already-processed part; document what they mean.
  \item A wrapper supplies the initial accumulator.
  \item Accumulative reverse is $\Theta(n)$ versus $\Theta(n^2)$ structural.
  \item Mutual recursion: functions that call each other, often mirroring states or data definitions.
\end{itemize}
